\section{Torsión integral, complejos y determinantes}
\label{sec:torsion-integral}

La Nota~\ref{sec:capa-fija} dejó una entrada de tamaño: acotar la
torsión de la homología de $\operatorname{GL}_N(\mathcal O_K)$
incluyendo todos los primos. Desarrollamos esa entrada. Primero
medimos la torsión de un complejo libre por los volúmenes de sus
fronteras. Después construimos un modelo finito para un subgrupo
de congruencia sin torsión y recuperamos el grupo completo mediante
un descenso integral de índice finito.

La construcción de Ash y el teorema de bases acotadas de
Bayer-Fluckiger, Emery y Houriet se reciben como resultados clásicos.
A partir de ellos probamos el recuento de celdas, la cota de las
fronteras, el descenso y la absorción cuantitativa. La conclusión
proporciona un exponente absoluto explícito para la entrada de
torsión de la Nota~\ref{sec:capa-fija}; las otras entradas de aquella
nota conservan su alcance.

\fuente{A93, pp.~72--74, localiza la estimación integral requerida.
Soulé~\cite{soule-torsion}, §§3.1--3.3, pp.~776--778, aporta el
antecedente geométrico y el umbral de primos que examinamos.
Bayer-Fluckiger--Emery--Houriet~\cite{bfeh}, Teorema~2.2, p.~74,
y Proposición~4.7, p.~80, proporciona las dos entradas geométricas.
Brown~\cite{brown-grupos}, Corolario~3.4, p.~22, y (5.6), p.~33,
y Weibel~\cite{weibel-homologia}, §§6.5 y~6.8, proporcionan la
transferencia y la sucesión espectral de una extensión.
No se incorpora un nuevo cotejo de declaraciones ni se ejecuta Lean.}

\subsection{Qué orden se debe controlar}

Sea $K$ un cuerpo de números, $n=[K:\Q]\ge2$,
$D=|\operatorname{disc}K|$, $N\ge5$ y
\[
 \Gamma=\operatorname{GL}_N(\mathcal O_K),\qquad
 T_{N,j}(K)=\#H_j(\Gamma,\Z)_{\mathrm{tors}}.
\]
La finitud de ese orden se deducirá de la construcción.
Todos los logaritmos sin subíndice son naturales.
Para un grupo abeliano finitamente generado $A$ escribimos
$t(A)=\#A_{\mathrm{tors}}$; en particular, $t(0)=1$.

La aplicación a la nota anterior usa $N=4e+6$ y
$0\le j\le2e+1$, con $e\ge1$. Obtendremos
\begin{equation}
 \log T_{4e+6,j}(K)
 \le \exp\!\left((2en)^{20}\log(2D)\right).
 \label{eq:ti-objetivo}
\end{equation}
El exponente es deliberadamente holgado. Su función es conservar
una dependencia polinómica en $e,n$ dentro del exponente y lineal
en $\log(2D)$. El orden de torsión de la izquierda es el completo.
No se localizan los coeficientes ni se elimina una parte primaria.

\subsection{La torsión homológica está en un conúcleo}

\begin{lema}[Frontera y torsión]
\label{lem:ti-conucleo}
Sea $C_\bullet$ un complejo de grupos abelianos libres de rango
finito, con $C_{-1}=0$. Para cada $j\ge0$ la inclusión de ciclos
induce un isomorfismo canónico
\begin{equation}
 H_j(C_\bullet)_{\mathrm{tors}}
 \simeq
 \bigl(C_j/\operatorname{im}d_{j+1}\bigr)_{\mathrm{tors}}.
 \label{eq:ti-conucleo}
\end{equation}
\end{lema}

\begin{proof}
La frontera $d_j$ induce la sucesión exacta
\[
 0\longrightarrow H_j(C_\bullet)
 \longrightarrow C_j/\operatorname{im}d_{j+1}
 \longrightarrow\operatorname{im}d_j\longrightarrow0.
\]
El último grupo es un subgrupo de $C_{j-1}$ y, por tanto, es
libre. Todo elemento de torsión del grupo central tiene imagen
cero, de modo que pertenece a $H_j(C_\bullet)$. La inclusión
conserva su orden y prueba la igualdad de las partes de torsión.
No hace falta elegir una escisión.
\end{proof}

Este paso evita estimar una base del núcleo de $d_j$.
Tampoco supone que $\operatorname{im}d_{j+1}$ sea un subretículo
primitivo: su falta de saturación es precisamente la torsión
que necesitamos medir.

\subsection{Saturación y determinantes}

Para $L\subset\Z^a$, su saturación es
$L^{\mathrm{sat}}=(L\otimes_\Z\Q)\cap\Z^a$.
Un elemento de $\Z^a/L$ es de torsión si y solo si algún
múltiplo entero no nulo pertenece a $L$. Por ello
\[
 (\Z^a/L)_{\mathrm{tors}}=L^{\mathrm{sat}}/L.
\]
Los covolúmenes siguientes se calculan en el subespacio real
generado por $L$, con la métrica euclídea de $\R^a$.

\begin{proposicion}[Cota de Gabber por columnas]
\label{prop:ti-gabber}
Sea $A:\Z^b\to\Z^a$ una matriz entera de rango $r$.
Si $v_1,\ldots,v_r$ son columnas linealmente independientes,
entonces
\begin{equation}
 t(\operatorname{coker}A)
 \le\prod_{i=1}^r\|v_i\|_2.
 \label{eq:ti-gabber}
\end{equation}
Si todas las columnas tienen norma a lo sumo $B\ge1$, resulta
\begin{equation}
 \log t(\operatorname{coker}A)
 \le r\log B\le\min(a,b)\log B.
 \label{eq:ti-gabber-log}
\end{equation}
Para $r=0$ el producto vacío vale uno.
\end{proposicion}

\begin{proof}
Pongamos $L=\operatorname{im}A$ y
$L_0=\sum_{i=1}^r\Z v_i$. Las inclusiones
$L_0\subset L\subset L^{\mathrm{sat}}$ tienen índice finito y
\[
 [L^{\mathrm{sat}}:L]
 \le[L^{\mathrm{sat}}:L_0]
 =\frac{\operatorname{covol}(L_0)}
        {\operatorname{covol}(L^{\mathrm{sat}})}.
\]
Si $w_1,\ldots,w_r$ es una base entera de $L^{\mathrm{sat}}$,
la fórmula de Cauchy--Binet expresa
$\det((\langle w_i,w_k\rangle)_{i,k})$ como la suma de los
cuadrados de sus menores de orden $r$. Es un entero positivo.
Así $\operatorname{covol}(L^{\mathrm{sat}})\ge1$.
Por Hadamard,
$\operatorname{covol}(L_0)\le\prod_i\|v_i\|_2$.
La identificación con la saturación prueba
\eqref{eq:ti-gabber}; tomar logaritmos da la otra desigualdad.
Si $r=0$, la imagen es cero y el conúcleo es libre.
\end{proof}

Los determinantes se toman en rango $r$, no necesariamente en
todo $\R^a$. El argumento admite un conúcleo con parte libre.
Es la forma de Gabber del Lema~3.1 de
Bayer-Fluckiger--Emery--Houriet; aquí hemos dado la prueba que
usaremos y conservado el rango efectivo.

\subsection{De las columnas a un complejo celular}

\begin{proposicion}[Tamaño y fronteras]
\label{prop:ti-complejo}
Sean $c_i=\operatorname{rank}C_i$ y $B_{j+1}\ge1$ una cota
para las normas de las columnas de $d_{j+1}$ en bases enteras.
Entonces
\begin{equation}
 \begin{aligned}
 \operatorname{rank}H_j(C_\bullet)&\le c_j,\\
 \log t(H_j(C_\bullet))
 &\le\min(c_j,c_{j+1})\log B_{j+1}.
 \end{aligned}
 \label{eq:ti-complejo}
\end{equation}
Si cada columna tiene a lo sumo $s$ entradas no nulas de
valor absoluto a lo sumo $a$, se puede tomar
$B_{j+1}=\max(1,a\sqrt{s})$.
\end{proposicion}

\begin{proof}
La cota de rango se obtiene del núcleo contenido en $C_j$.
El Lema~\ref{lem:ti-conucleo} reduce la torsión al conúcleo
de $d_{j+1}$ y la Proposición~\ref{prop:ti-gabber} la acota.
La última afirmación usa que la suma de los cuadrados de
las entradas de cada columna es a lo sumo $sa^2$.
\end{proof}

Para celdas poliedrales con a lo sumo $b$ facetas geométricas,
cada incidencia tiene signo $\pm1$ antes de pasar al cociente.
Varias facetas pueden identificarse después. La estimación
siempre válida para la columna resultante es
\begin{equation}
 \|d\sigma\|_2\le\|d\sigma\|_1\le b.
 \label{eq:ti-incidencias}
\end{equation}
Una mejora a $\sqrt b$ exigiría controlar esas identificaciones.
Usaremos \eqref{eq:ti-incidencias}, suficiente para la cota final.

El número de celdas por sí solo no acota la torsión: el complejo
$\Z\xrightarrow{\times a}\Z$ tiene un generador en cada uno
de dos grados y conúcleo $\Z/a$. La frontera contiene el dato
de tamaño. Al tensorizar con $\Q$ ese conúcleo desaparece;
el complejo racional tampoco permite recuperar su orden.

\subsection{Homología de un grupo finito con coeficientes}

Prepararemos el descenso sin suponer trivial la acción en los
coeficientes. Sea $F$ finito, $m=|F|\ge2$, y sea $M$ un
$\Z[F]$-módulo finitamente generado como grupo abeliano.
Escribamos $r=\operatorname{rank}M$ y $t=t(M)$.

\begin{proposicion}[Control de todos los grados]
\label{prop:ti-grupo-finito}
Se tienen
\begin{equation}
 \begin{aligned}
 t(H_0(F,M))&\le m^r t,\\
 \#H_p(F,M)&\le(m^r t)^{(m-1)^p}\qquad(p>0).
 \end{aligned}
 \label{eq:ti-grupo-finito}
\end{equation}
En particular, los grupos de grado positivo son finitos.
\end{proposicion}

\begin{proof}
Sea $T=M_{\mathrm{tors}}$ y $A=M/T$, libre de rango $r$.
Los coinvariantes dan una sucesión exacta por la derecha
$T_F\to M_F\to A_F\to0$. Su primer término es finito
de orden a lo sumo $t$. La norma induce
$\mathcal N:A_F\to A^F$, por
$[a]\mapsto\sum_{g\in F}ga$.
La composición con $A^F\to A_F$ es multiplicar por $m$.
Como $A^F$ es libre, la torsión de $A_F$ está anulada por $m$.
El grupo $A_F$ tiene a lo sumo $r$ generadores, de modo que
$t(A_F)\le m^r$. La torsión de $M_F$ tiene núcleo de orden
a lo sumo $t$ e imagen contenida en la torsión de $A_F$.
Esto prueba la primera desigualdad.

Para $p>0$ recibimos la identidad clásica de transferencia:
$mH_p(F,M)=0$, válida para cualquier acción en $M$.
Se deduce al restringir al subgrupo trivial y volver a incluir:
la composición es multiplicar por el índice, y la homología
positiva del grupo trivial es cero.

El complejo bar normalizado tiene grupo de cadenas
$M^{(m-1)^p}$ en grado $p$: se suprimen los símbolos que
contienen la identidad. Su grupo de ciclos $Z_p$ tiene rango
a lo sumo $r(m-1)^p$ y torsión de orden a lo sumo
$t^{(m-1)^p}$. Como $m$ anula su cociente $H_p(F,M)$,
el mapa de ciclos factoriza por $Z_p/mZ_p$. Para todo grupo
abeliano finitamente generado $Z$,
\[
 \#(Z/mZ)=m^{\operatorname{rank}Z}\#(Z_{\mathrm{tors}}/mZ_{\mathrm{tors}})
 \le m^{\operatorname{rank}Z}t(Z).
\]
Aplicar esta fórmula a $Z_p$ prueba la segunda desigualdad
y, a la vez, la finitud.
\end{proof}

El factor $m^r$ mide la torsión que una acción puede crear al
pasar a coinvariantes. Por ejemplo, si $F=C_2$ actúa en $\Z$
por cambio de signo, sus coinvariantes son $\Z/2$ aunque
el módulo original sea libre.

\clearpage
\subsection{Descenso integral por índice finito}

\begin{teorema}[Una extensión de grupos]
\label{thm:ti-descenso}
Sea $1\to\Gamma'\to\Gamma\to F\to1$ exacta, con
$m=|F|\ge2$, y supongamos que $H_q(\Gamma',\Z)$ es
finitamente generado para $0\le q\le j$.
Pongamos
$r_q=\operatorname{rank}H_q(\Gamma',\Z)$ y
$t_q=t(H_q(\Gamma',\Z))$. Entonces $H_j(\Gamma,\Z)$ es
finitamente generado y
\begin{equation}
 t(H_j(\Gamma,\Z))
 \le\prod_{q=0}^j(m^{r_q}t_q)^{(m-1)^{j-q}}.
 \label{eq:ti-descenso}
\end{equation}
Equivalentemente,
\begin{equation}
 \log t(H_j(\Gamma,\Z))
 \le\sum_{q=0}^j(m-1)^{j-q}
          \bigl(r_q\log m+\log t_q\bigr).
 \label{eq:ti-descenso-log}
\end{equation}
Si $m=1$, se tiene simplemente $\Gamma'=\Gamma$.
\end{teorema}

\begin{proof}
Usamos la sucesión espectral homológica de
Lyndon--Hochschild--Serre
\begin{equation}
 E^2_{p,q}=H_p(F,H_q(\Gamma',\Z))
 \ \Longrightarrow\ H_{p+q}(\Gamma,\Z).
 \label{eq:ti-lhs}
\end{equation}
La acción de $F$ procede de conjugación; no necesitamos que
sea trivial. Esta sucesión de primer cuadrante proporciona
en grado total $j$ una filtración finita con cocientes
$E^\infty_{p,j-p}$.

Para $p>0$ y $0\le q\le j$, los términos $E^2_{p,q}$ son finitos por la
Proposición~\ref{prop:ti-grupo-finito}. Sus términos posteriores
son subcocientes y su orden no aumenta.
El término $E^2_{0,j}$ puede tener parte libre. Debemos
tratarlo aparte: un cociente arbitrario de un grupo libre
puede crear torsión.

Los diferenciales homológicos tienen bidegrado $(-a,a-1)$
en la página $a\ge2$. Desde $(0,j)$ no sale ninguno.
Solo llegan imágenes de $E^a_{a,j-a+1}$, que son finitas
porque su primera coordenada es positiva. En cada paso
se divide, por tanto, por un subgrupo finito contenido en
la torsión. Para $B$ finitamente generado y $U\subset B$
finito, $(B/U)_{\mathrm{tors}}=B_{\mathrm{tors}}/U$.
Resulta
$t(E^\infty_{0,j})\le t(E^2_{0,j})$.

Todos los cocientes de la filtración son finitamente
generados, luego también lo es $H_j(\Gamma,\Z)$.
En una sucesión exacta $0\to A\to B\to C\to0$ de tales
grupos se cumple $t(B)\le t(A)t(C)$: el núcleo de la
restricción a torsión es $A_{\mathrm{tors}}$ y su imagen
está contenida en $C_{\mathrm{tors}}$.
Aplicando esto a la filtración y usando
\eqref{eq:ti-grupo-finito}, con $p=j-q$, obtenemos
\eqref{eq:ti-descenso}. Tomar logaritmos da
\eqref{eq:ti-descenso-log}.
\end{proof}

El descenso utiliza los rangos de la cubierta además de sus
torsiones. No se intenta invertir el índice $m$: sus factores
primos quedan incluidos en el coste.

\subsection{Un subgrupo de congruencia sin torsión}

Definimos
\[
 \Gamma(3)=\ker\!\left(
 \operatorname{GL}_N(\mathcal O_K)\longrightarrow
 \operatorname{GL}_N(\mathcal O_K/3\mathcal O_K)\right).
\]
El anillo residual puede tener nilpotentes; el argumento
no exige que $3$ sea no ramificado.

\begin{lema}[Nivel tres]
\label{lem:ti-nivel-tres}
El grupo $\Gamma(3)$ es normal, sin torsión y de índice
\begin{equation}
 m=[\Gamma:\Gamma(3)]\le3^{nN^2}.
 \label{eq:ti-indice}
\end{equation}
\end{lema}

\begin{proof}
La normalidad es inmediata. Como $\mathcal O_K$ es libre
de rango $n$ sobre $\Z$, el anillo $\mathcal O_K/3\mathcal O_K$
tiene $3^n$ elementos. El número de todas las matrices
$N\times N$ sobre él es $3^{nN^2}$, lo que acota la imagen.

Eligiendo una base entera de $\mathcal O_K$, la restricción
de escalares da una inclusión
$\Gamma\hookrightarrow\operatorname{GL}_{nN}(\Z)$ y lleva
$\Gamma(3)$ al núcleo de reducción módulo $3$.
Si ese núcleo tuviera un elemento de orden finito no trivial,
alguna potencia tendría orden primo $p$. Escribámosla
$g=I+3^aB$, con $a\ge1$ y alguna entrada de $B$ no divisible
por $3$.

Si $p\ne3$, dividir $g^p-I$ por $3^a$ y reducir módulo $3$
da $pB\ne0$, contradicción. Si $p=3$, la identidad exacta es
\[
 \frac{g^3-I}{3^{a+1}}
   =B+3^aB^2+3^{2a-1}B^3.
\]
Su reducción módulo $3$ es $B\ne0$, otra contradicción.
Así no existe torsión no trivial.
\end{proof}

\subsection{Las dos entradas geométricas}

Fijemos $\Lambda=\mathcal O_K^N$ y sea $\Sigma$ el conjunto
de las $n$ inmersiones de $K$ en $\C$.
Una métrica hermítica $h$ da la forma euclídea real
\[
 q_h(x)=\sum_{\sigma\in\Sigma}h_\sigma(x_\sigma,x_\sigma).
\]
Las dos inmersiones de cada par complejo se cuentan en esta
suma. Denotamos por $M(h)$ el conjunto de vectores no nulos
de $\Lambda$ con $q_h(x)=1$.

\paragraph{Retracto de Ash, recibido.}
El espacio $\widetilde W$ de métricas de mínimo uno para las
que $M(h)$ genera $K^N$ sobre $K$ es contráctil y tiene
una estructura celular poliedral, localmente finita,
compatible con $\Gamma$. Una celda abierta está determinada
por su conjunto exacto de vectores mínimos. Sus caras se
obtienen añadiendo igualdades $q_h(x)=1$; las clausuras
son poliedros con incidencias de facetas $\pm1$.
Los estabilizadores de celdas son finitos y el cociente
por $\Gamma$ es compacto.
Esta es la construcción descrita en
Bayer-Fluckiger--Emery--Houriet, §2.3, Teorema~2.2,
y Soulé, §3.1.

El espacio real de todas las formas tiene dimensión
\begin{equation}
 d=r_1\frac{N(N+1)}2+r_2N^2
 \le \delta:=n\frac{N(N+1)}2\le nN^2.
 \label{eq:ti-dimension}
\end{equation}
Usaremos $\delta$ como sobrecota; no necesitamos conocer
la dimensión exacta del retracto.

\paragraph{Bases acotadas, recibido.}
La Proposición~4.7 de Bayer-Fluckiger--Emery--Houriet,
para $n\ge2$ y $N\ge5$, proporciona en cada retículo
$(\Lambda,h)$ de $\widetilde W$ una base sobre
$\mathcal O_K$ cuyos vectores tienen norma a lo sumo
\begin{equation}
 B=4N^2 n^{5Nn^2}D^{6N(n+1)}.
 \label{eq:ti-base-acotada}
\end{equation}
Hemos aumentado su constante suprimiendo el factor $2^{-N}$.
Desde esta entrada construiremos el conjunto finito de vectores
necesario para contar. No usaremos la comparación homológica
fuera de pequeños primos de su Corolario~2.3.

\subsection{Un conjunto finito de vectores mínimos}

\begin{proposicion}[Coordenadas mediante volumen y empaquetamiento]
\label{prop:ti-phi}
Pongamos
\begin{equation}
 s=nN,\qquad
 R=n\,s^s B^{2(s-1)}D^N,\qquad
 P=(3\sqrt R)^s.
 \label{eq:ti-R-P}
\end{equation}
Existe un conjunto finito $\Phi\subset\Lambda$ con
$\#\Phi\le P$ tal que para cada $h\in\widetilde W$ algún
$\gamma\in\Gamma$ satisface $\gamma M(h)\subset\Phi$.
\end{proposicion}

\begin{proof}
Tomemos una base acotada $e_1,\ldots,e_N$ y escribamos
$H_\sigma=(h_\sigma(e_i,e_k))_{i,k}$.
Para $x=\sum_i x_i e_i\in M(h)$,
Cauchy--Schwarz para la forma $H_\sigma$ y Hadamard
en su menor principal dan
\[
 |\sigma(x_i)|^2
 \le (H_\sigma^{-1})_{ii}\,h_\sigma(x)
 \le \det(H_\sigma)^{-1}B^{2(N-1)}.
\]
Aquí $h_\sigma(x)\le q_h(x)=1$.
También $\det H_\sigma\le B^{2N}$, de modo que
\begin{equation}
 \sum_{\sigma\in\Sigma}|\sigma(x_i)|^2
 \le n B^{2(s-1)}
           \prod_{\sigma\in\Sigma}\det(H_\sigma)^{-1}.
 \label{eq:ti-coordenadas}
\end{equation}

El retículo euclídeo $(\Lambda,q_h)$ tiene dimensión $s$,
mínimo uno y covolumen $v$. Las bolas abiertas de radio
$1/2$ se inyectan en el toro cociente. Una de ellas contiene
un cubo de lado $1/\sqrt{s}$, por lo que
$v\ge s^{-s/2}$.
La fórmula del discriminante, con las dos inmersiones
complejas incluidas en la métrica, da
\[
 v^2=D^N\prod_{\sigma\in\Sigma}\det H_\sigma.
\]
En efecto, para la métrica de traza estándar el covolumen de
$\mathcal O_K$ es $\sqrt D$; cambiar a las matrices
$H_\sigma$ multiplica el cuadrado del volumen por ese producto.
Por tanto
$\prod_\sigma\det(H_\sigma)^{-1}\le s^sD^N$.
La desigualdad \eqref{eq:ti-coordenadas} queda acotada por $R$.

En una base fija $f_1,\ldots,f_N$ de $\Lambda$ definimos
\[
 \Phi=\left\{\sum_i a_i f_i:
       a_i\in\mathcal O_K,\ 
       \sum_{\sigma\in\Sigma}|\sigma(a_i)|^2\le R
       \text{ para cada }i\right\}.
\]
El automorfismo que lleva cada $e_i$ a $f_i$ transporta
$M(h)$ dentro de $\Phi$.

Para contar, la distancia entre dos enteros algebraicos
distintos en la métrica de traza es al menos uno.
En efecto, para $a\ne0$ el entero $|N_{K/\Q}(a)|\ge1$,
y la desigualdad aritmético-geométrica da
$\sum_\sigma|\sigma(a)|^2\ge n$.
Las bolas de radio $1/2$ alrededor de los puntos de norma
a lo sumo $\sqrt R$ son disjuntas y están dentro de una
bola de radio $\sqrt R+1/2$. Hay a lo sumo
$(1+2\sqrt R)^n\le(3\sqrt R)^n$ puntos, porque $R\ge1$.
Las $N$ coordenadas se eligen por separado, luego
$\#\Phi\le(3\sqrt R)^{nN}=P$.
\end{proof}

La identidad de covolúmenes conserva la normalización de
traza en los lugares complejos. La estimación inferior del
volumen usa solo el mínimo del retículo; por ello el cálculo
no necesita una elección adicional de unidades.

\subsection{Contar ecuaciones independientes y facetas}

\begin{proposicion}[Complejidad celular]
\label{prop:ti-celdas}
Si $\Phi$ tiene la propiedad de la
Proposición~\ref{prop:ti-phi}, entonces hay a lo sumo $P^\delta$
órbitas de celdas de cada dimensión bajo $\Gamma$.
Cada celda de $\widetilde W$ tiene a lo sumo
\begin{equation}
 \beta=P^{N+1}
 \label{eq:ti-facetas}
\end{equation}
facetas geométricas.
\end{proposicion}

\begin{proof}
En el espacio real de formas, $L_x(h)=q_h(x)$ es lineal
en $h$. Una celda de dimensión $q$ y conjunto de mínimos $M$
tiene envolvente afín dada por las igualdades $L_x(h)=1$
para $x\in M$. Para justificarlo, cerca de una métrica
$h_0$ de la celda se tiene $q_h\ge q_{h_0}/2$.
Solo los finitos vectores con $q_{h_0}(x)\le4$ pueden
amenazar las desigualdades estrictas. Los que no están
en $M$ tienen un margen positivo común. Así la celda
contiene una vecindad relativa de $h_0$ en el espacio
afín de sus igualdades. Podemos escoger $d-q$ ecuaciones
independientes que definan esa misma envolvente.

La envolvente determina la celda: $M$ es exactamente el
conjunto de vectores de $\Lambda\setminus\{0\}$ cuya ecuación
vale idénticamente en ella. Una igualdad que valga en el
interior relativo vale en toda la envolvente; una igualdad
que valga en la envolvente vale, en particular, en cada
punto de la celda. Las demás desigualdades son estrictas.
No se supone que $|M|=d-q$; puede haber muchas ecuaciones
redundantes.

Transportando $M$ dentro de $\Phi$, el número de posibilidades
es a lo sumo $\binom{\#\Phi}{d-q}\le P^\delta$.
Un subconjunto independiente determina a lo sumo una celda,
lo que basta para acotar las órbitas.

Para las facetas fijemos una celda con mínimos $M$ y una
base $v_1,\ldots,v_N$ de $K^N$ contenida en $M$.
Cada faceta añade al menos una ecuación $L_x=1$ independiente
de las anteriores; una de ellas determina su envolvente
afín dentro de la de la celda. No hace falta que su conjunto
completo de mínimos sea $M\cup\{x\}$.

Sea $M'$ el conjunto de mínimos de la faceta.
Elegimos $\gamma$ con $\gamma M'\subset\Phi$.
El mapa $\gamma$ queda determinado por las $N$ imágenes
$\gamma v_i\in\Phi$; una vez fijado, $x$ queda determinado
por $\gamma x\in\Phi$. Hay a lo sumo
$(\#\Phi)^{N+1}\le P^{N+1}$ posibilidades.
Cada una determina a lo sumo una faceta. Esta cuenta se
realiza antes de tomar el cociente por el grupo.
\end{proof}

La elección de ecuaciones independientes hace explícito el
recuento sugerido por Soulé, Proposición~3, y
Bayer-Fluckiger--Emery--Houriet, §3.3. Los vectores mínimos
y sus ecuaciones pueden tener dependencias; el argumento
solo usa el rango del sistema afín.

\subsection{Un complejo libre y la torsión del grupo completo}

\begin{proposicion}[Modelo integral de la cubierta]
\label{prop:ti-cubierta}
El cociente $X'=\Gamma(3)\backslash\widetilde W$ es un
modelo celular finito de $B\Gamma(3)$.
Sus cadenas celulares enteras son libres y satisfacen
\begin{equation}
 c_q\le mP^\delta,\qquad
 \|d_{q+1}\text{ en cada columna}\|_2\le P^{N+1}.
 \label{eq:ti-cubierta}
\end{equation}
En consecuencia, para todo $q\ge0$,
\begin{equation}
 r_q\le mP^\delta,\qquad
 \log t_q\le mP^\delta(N+1)\log P.
 \label{eq:ti-cubierta-homologia}
\end{equation}
\end{proposicion}

\begin{proof}
Un elemento que estabiliza una celda de $\widetilde W$
pertenece a su estabilizador finito. Su intersección con
$\Gamma(3)$ es trivial por el Lema~\ref{lem:ti-nivel-tres}.
Por tanto el subgrupo actúa libremente sobre las celdas
y sus puntos. La acción es propia: $\Gamma$ es discreto
en $\operatorname{GL}_N(K\otimes_\Q\R)$, y el espacio
de formas positivas es un espacio homogéneo de ese grupo
real con estabilizador compacto. La contractibilidad de
$\widetilde W$ hace del cociente un modelo de
$B\Gamma(3)$.

Cada órbita de $\Gamma$ se descompone en a lo sumo $m$
órbitas de $\Gamma(3)$. La Proposición~\ref{prop:ti-celdas}
acota así el número de celdas de cada dimensión; hay
solo un número finito de dimensiones. Elegir una
orientación por celda del cociente da bases enteras.

Aunque el cociente pueda identificar caras distintas de
una celda, su frontera suma a lo sumo $P^{N+1}$
incidencias con signo. La desigualdad
\eqref{eq:ti-incidencias} prueba la segunda cota.
La Proposición~\ref{prop:ti-complejo} da
\eqref{eq:ti-cubierta-homologia}. No se necesita subdividir
ni invertir los órdenes de estabilizadores.
\end{proof}

\begin{teorema}[Cota integral aritmética]
\label{thm:ti-aritmetica}
Bajo las dos entradas geométricas citadas, para todo
cuerpo $K$ de grado $n\ge2$, todo $N\ge5$ y todo $j\ge0$,
$H_j(\operatorname{GL}_N(\mathcal O_K),\Z)$ es
finitamente generado y
\begin{equation}
 \log T_{N,j}(K)
 \le (j+1)m^{j+1}P^\delta
       \bigl(\log m+(N+1)\log P\bigr),
 \label{eq:ti-aritmetica}
\end{equation}
con $m,\delta,B,R,P$ como arriba.
\end{teorema}

\begin{proof}
Para $m\ge2$ aplicamos el Teorema~\ref{thm:ti-descenso}
al cociente finito $\Gamma/\Gamma(3)$.
Cada sumando de \eqref{eq:ti-descenso-log} es a lo sumo
\[
 m^{j-q}\,mP^\delta
       \bigl(\log m+(N+1)\log P\bigr).
\]
Hay $j+1$ sumandos y $m^{j-q}\le m^j$.
Esto prueba \eqref{eq:ti-aritmetica}.
Si $m=1$, la Proposición~\ref{prop:ti-cubierta} da directamente
$\log T_{N,j}(K)\le P^\delta(N+1)\log P$, menor o igual
que la derecha de \eqref{eq:ti-aritmetica}.
\end{proof}

El complejo celular finito calcula la homología de
$\Gamma(3)$. La sucesión espectral calcula el coste de
recuperar $\Gamma$, cuya torsión impide usar directamente
un modelo finito libre del mismo tipo.

\subsection{Un exponente absoluto para la capa fija}

\begin{proposicion}[Absorción explícita]
\label{prop:ti-exponente}
Para $e\ge1$, $n\ge2$, $N=4e+6$ y $0\le j\le2e+1$,
la cota \eqref{eq:ti-aritmetica} implica
\eqref{eq:ti-objetivo}. Se puede tomar $C_T=20$ en
\eqref{eq:cf-torsion-entrada}.
\end{proposicion}

\begin{proof}
Pongamos $u=2en\ge4$ y $Y=\log(2D)\ge1/2$.
Entonces
\begin{equation}
 n\le u,\quad N\le u^2,\quad j+1\le u,\quad
 s\le u^3,\quad \delta\le u^5,\quad \log m\le2u^5.
 \label{eq:ti-parametros}
\end{equation}
Por ejemplo, $N\le10e\le(2en)^2$;
$j+1\le2e+2\le4e\le2en$.
Usando $\log u\le u$ y $\log D\le Y$ en
\eqref{eq:ti-base-acotada}, obtenemos
\[
 \log B
 \le2u^2+5u^5+12u^3Y
 \le(4u^2+10u^5+12u^3)Y
 \le u^7Y.
\]
La última desigualdad vale para $u\ge4$, como se ve
dividiendo entre $u^7$. De \eqref{eq:ti-R-P} se sigue
\[
 \begin{aligned}
 \log R
 &=\log n+s\log s+2(s-1)\log B+N\log D\\
 &\le u+3u^4+2u^{10}Y+u^2Y
 \le u^{11}Y,\\
 \log P
 &=s\bigl(\log3+\tfrac12\log R\bigr)
 \le u^3\bigl(2+\tfrac12u^{11}Y\bigr)
 \le u^{14}Y.
 \end{aligned}
\]
Todas las constantes se absorbieron usando $u\ge4$ y
$Y\ge1/2$, sin acotar el discriminante por una constante.

Escribamos $Q_0=\log m+(N+1)\log P>0$.
Las cotas anteriores dan $Q_0\le4u^{16}Y$.
El logaritmo de la derecha de \eqref{eq:ti-aritmetica}
está acotado por
\[
 \begin{aligned}
 \log(j+1)+(j+1)\log m+\delta\log P+\log Q_0
 &\le u+2u^6+u^{19}Y+4u^{16}Y\\
 &\le(6u^6+u^{19}+4u^{16})Y\\
 &\le2u^{19}Y\le u^{20}Y.
 \end{aligned}
\]
Usamos $\log Q_0\le Q_0$ y
$6u^6+4u^{16}\le u^{19}$ para $u\ge4$.
Al exponenciar queda \eqref{eq:ti-objetivo}.
\end{proof}

En la Nota~\ref{sec:capa-fija}, la fórmula
\eqref{eq:cf-exponente} da ahora
\begin{equation}
 C_K=\max(5\cdot20+9,29)+2=111.
 \label{eq:ti-CK}
\end{equation}
Esta sustitución descarga la entrada de torsión.
La deducción de la cota de $H^2$ sigue recibiendo el rango
filtrado, la estabilidad integral, los ciclos, las comparaciones
de clases y Selmer y la extensión auxiliar que allí se enumeran.

\subsection{Qué ocurre con los pequeños primos}

El descenso por $\Gamma(3)$ evita decidir qué primos se pueden
suprimir. Conviene, de todos modos, precisar el antecedente
citado en la nota anterior.

\begin{lema}[Un umbral válido por restricción de escalares]
\label{lem:ti-primos}
Para este lema admitimos cualquier $N\ge1$.
Si $\gamma\in\operatorname{GL}_N(\mathcal O_K)$ tiene orden
primo $p$, entonces
\begin{equation}
 p\le nN+1.
 \label{eq:ti-primos}
\end{equation}
En consecuencia, todo primo que divide el orden de un
subgrupo finito de $\Gamma$ satisface esa desigualdad.
\end{lema}

\begin{proof}
La restricción de escalares representa a $\gamma$ como
una matriz racional de tamaño $nN$ y del mismo orden.
Su polinomio mínimo divide $X^p-1$ y tiene algún factor
distinto de $X-1$. Como
$\Phi_p(X)=1+X+\cdots+X^{p-1}$ es irreducible sobre $\Q$,
ese factor tiene grado $p-1$. Así $p-1\le nN$.
La afirmación sobre subgrupos finitos sigue del teorema de
Cauchy, que proporciona un elemento de orden $p$.
\end{proof}

El Lema~9 de Soulé (2003), pp.~777--778, imprime el
umbral $p\le1+\max(n,N)$. Tal enunciado, en esa generalidad,
admite un contraejemplo. Tomemos
\begin{equation}
 K=\Q(\sqrt5),\qquad
 t=\frac{\sqrt5-1}{2},\qquad
 \gamma=\begin{pmatrix}0&-1\\1&t\end{pmatrix}
       \in\operatorname{SL}_2(\mathcal O_K).
 \label{eq:ti-contraejemplo}
\end{equation}
Tenemos $t^2+t=1$, $\gamma^2=t\gamma-I$ y, sucesivamente,
\[
 \gamma^3=-t\gamma-tI,\qquad
 \gamma^4=-\gamma+tI,\qquad
 \gamma^5=I.
\]
Como $\gamma\ne I$ y $5$ es primo, su orden es $5$.
Aquí $n=N=2$ y $1+\max(n,N)=3$.
La cota \eqref{eq:ti-primos}, en cambio, da $5$.

El paso problemático de la prueba impresa es la afirmación
de que $\Phi_p$ permanece irreducible sobre $K$ cuando
$K$ no contiene todas las raíces $p$-ésimas de la unidad.
Un subcuerpo propio del cuerpo ciclotómico puede reducir
su grado de irreducibilidad. El ejemplo anterior lo realiza.
Esta observación corrige ese umbral concreto; no se utiliza
para atribuir un fallo a todos los resultados del artículo.
Nuestra estimación integral se apoya en el complejo de la
cubierta y en el descenso demostrado, y no usa ese lema.

\clearpage
\subsection{Resultado obtenido y entradas que permanecen}

\begin{center}
\small
\begin{tabularx}{\textwidth}{@{}p{0.28\textwidth}Y@{}}
\toprule
Paso & Alcance en esta nota\\
\midrule
Complejos y determinantes &
Pruebas completas de saturación, Gabber, normas de frontera
y torsión homológica.\\
Descenso de índice finito &
Prueba de la cota integral desde transferencia y la sucesión
espectral clásicas, incluidas acciones no triviales.\\
Geometría aritmética &
Retracto de Ash y bases acotadas recibidos; construcción de
$\Phi$, recuento, facetas y cubierta desarrollados aquí.\\
Torsión y capa fija &
$C_T=20$ demostrado desde esas entradas geométricas;
$C_K=111$ bajo las demás entradas de la Nota~\ref{sec:capa-fija}.\\
Continuación &
Rango filtrado, estabilidad y realización de ciclos;
comparaciones de clases, unidades y Selmer.\\
\bottomrule
\end{tabularx}
\end{center}

La cota uniforme cubre cuerpos reales o imaginarios de grado
al menos dos. Para la aplicación a $H^2$ siguen vigentes
las hipótesis de cuerpo totalmente imaginario y raíces de
la unidad de la Nota~\ref{sec:capa-fija}.
No se ha probado aquí una cota integral para los grupos de
teoría $K$ mediante Hurewicz; la cantidad estimada es la
torsión de la homología del grupo lineal.

El siguiente bloque puede concentrarse en el rango de la
imagen homológica y la realización de los ciclos de A93.
Esas entradas deben producir los mismos mapas y condiciones
que usa la cancelación de la altura. La nueva cota de tamaño
no construye por sí sola tales mapas.

\clearpage
\subsection{Ejercicios con soluciones}

\begin{ejercicio}[Una imagen no saturada]
Sea $A:\Z^2\to\Z^3$ la matriz con columnas
$(2,0,0)$ y $(0,3,0)$. Calcular el conúcleo, su torsión
y la cota de Gabber. ¿Qué cambia al pasar a $\Q$?
\end{ejercicio}

\paragraph{Solución.}
Se tiene
$\operatorname{coker}A\simeq\Z/2\oplus\Z/3\oplus\Z$.
La saturación de la imagen es
$\Z(1,0,0)\oplus\Z(0,1,0)$ y su índice sobre la imagen
es $6$. Las columnas son ortogonales, de normas $2$ y $3$;
Gabber da exactamente $6$. Sobre $\Q$ queda solo un
conúcleo de dimensión uno, sin información sobre esos factores.

\begin{ejercicio}[Facetas identificadas]
Un cociente celular identifica $b$ facetas de una celda
con una misma celda de dimensión menor y con la misma
orientación. Comparar las normas $\ell^1$ y $\ell^2$
de la columna de frontera resultante.
\end{ejercicio}

\paragraph{Solución.}
La columna tiene una única entrada igual a $b$.
Ambas normas valen $b$, y para $b>1$ exceden $\sqrt b$.
Contar facetas antes de tomar el cociente sí permite
$\|d\sigma\|_2\le\|d\sigma\|_1\le b$.
La coincidencia de varias caras explica por qué conservamos
esa forma del presupuesto.

\begin{ejercicio}[Rangos en el descenso]
Sea $F=C_2$ actuando en $M=\Z$ por $x\mapsto-x$.
Calcular $H_0(F,M)$ y explicar por qué una cota de
descenso que solo utilice $t(M)$ sería insuficiente.
\end{ejercicio}

\paragraph{Solución.}
La relación de coinvariantes es $-x-x=-2x$, luego
$H_0(F,M)=\Z/2$. Aunque $t(M)=1$, los coinvariantes
tienen torsión de orden $2$.
La Proposición~\ref{prop:ti-grupo-finito} incorpora el coste
$m^{\operatorname{rank}M}=2$ y alcanza el orden exacto.
En la sucesión espectral, el rango de la homología de la
cubierta participa por esta misma razón.

\begin{ejercicio}[Dos cocientes distintos]
Comparar el efecto sobre la torsión de dividir $\Z$
por $2\Z$ y de dividir $\Z\oplus\Z/6$ por su subgrupo
$0\oplus\langle\bar3\rangle$.
¿Cuál corresponde al paso por $E^a_{0,j}$?
\end{ejercicio}

\paragraph{Solución.}
El primer cociente es $\Z/2$: un subgrupo infinito puede
crear torsión al tomar el cociente.
El segundo es $\Z\oplus\Z/3$: dividir por un subgrupo
finito reduce la torsión de orden $6$ a orden $3$.
En $E^a_{0,j}$ solo entran imágenes de grupos finitos;
por eso se aplica el segundo mecanismo. No sería correcto
acotar la torsión de cualquier subcociente por la torsión
del grupo inicial.

\begin{ejercicio}[Hipótesis y exponentes]
Tomar $e=1$ y un cuerpo totalmente imaginario de grado
$n=2$ y discriminante absoluto $D$.
Indicar $N$, los grados homológicos cubiertos y la cota
obtenida. Explicar qué falta para usar $C_K=111$.
\end{ejercicio}

\paragraph{Solución.}
Se tiene $N=10$, $u=2en=4$ y $0\le j\le3$.
La conclusión es
\[
 \log T_{10,j}(K)\le\exp\!\left(4^{20}\log(2D)\right)
                   =(2D)^{4^{20}}.
\]
Esta cota procede de Ash, de las bases acotadas y de las
deducciones de esta nota. Para obtener la cota de $H^2$
con $C_K=111$ hay que verificar además las hipótesis sobre
raíces de la unidad y todas las entradas restantes de la
Nota~\ref{sec:capa-fija}, incluida la extensión auxiliar
cuando $2e+1<r_2(K)$ no se cumple. En este ejemplo
$r_2(K)=1$, por lo que esa desigualdad falla.
El exponente explícito no reemplaza ninguna de esas condiciones.
